\par\SourceLeaf{165}
If, however, the changes are for the present to be kept to the absolute minimum permissible, the commission considers that this can be done in the gentlest possible way by the following plan, which it therefore conditionally recommends for adoption.

I.\quad The annual contributions are raised to 3 Louis d'or.

II.\quad The pensions remain for the present at 250 Rthl., as long as the number of widows\footnote{For brevity, only widows are mentioned here throughout; it is, however, self-evident that an orphan's pension must always be counted as a widow's pension.}, excluding Professor G.'s widow, does not exceed 18, the present number. Otherwise the pension is determined by adding to a fixed part of 200 Rthl. the dividend falling to each widow when 900 Rthl. (as 18 times 50 Rthl.) is shared equally among those present. At the highest number of widows that has yet occurred, namely 21 (excluding Prof. G.'s widow), the pension would thus still be 242 Rthl. 20 Ggr. 7 Pf. in gold.

It is readily seen that this arrangement does not yet cover the deficit in full: the calculation in fact shows that, assuming it takes effect from 1 October 1845, 2963 Rthl. still remain uncovered, and proportionately more if the higher contribution begins later. If, however, the higher interest income still enjoyed at present continues for the time being without substantial diminution, and no other considerable losses occur, one might hope that this remaining deficit will be covered of itself after a moderate number of years. But, if it adopts this plan, the association must not conceal from itself that, with its resources so narrowly matched to its obligations, the widows' fund resembles a heavily laden ship sailing in shallow waters, whose safety lies only in continually and vigilantly taking soundings.

The commission therefore considers it absolutely necessary that the above regulation be accompanied by

III.\quad a careful new balance calculation, to be repeated at suitable intervals (say every 10 or 5 years), with further assistance reserved as necessary according to its results. As a necessary preparatory measure, precise records of the personal circumstances of all participants relevant to the widows' fund would have to be introduced from now on, as already indicated at several points in the memorandum.

The commission does not consider it necessary to develop the advantages of this regulation at length here, and draws attention only to the following. If the first or second revision shows that the balance has not merely improved, but has improved far more than could have been hoped, it will be permissible to raise the pensions further from 250 Rthl. to 260 Rthl. when the number of widows is below 18, or, when the number is high, to raise the whole supplementary dividend from 900 Rthl. to 1080 Rthl. (that is, 18 times 60 Rthl.). If, on the other hand, the revision shows a deterioration in the balance, one will have to accept some further exertion, but will also have all the more reason to congratulate oneself on not having sat with one's hands in one's lap in 1846. In either case, however, further measures can be taken with an awareness of their safety or their necessity.

It need not be pointed out that adoption of this regulation automatically removes the beginning of §. 10 of the regulations, whose unclear wording must be regarded as a main source of the present trouble.

In applying the rules contained in this plan, the widows' fund must make no
\SourceLeaf{166}
distinction between the present widows and those who join their number in future, for only on the assumption of entirely equal treatment has it been possible to reduce the deficit to the moderate remainder of 2963 Rthl. Moreover, if the earlier widows were granted an exceptional preference \textit{from the fund}, the new widows would consider themselves wronged, since the possibility of limiting the modifications of the pensions to so mild a level is due solely to the increase in contributions, which the husbands of the new widows will have paid, but those of the earlier widows did not. Since, however, it also seems desirable to forestall any dissatisfaction among the present widows, even if not founded on a right, the commission proposes
\begin{quote}
that the university address a request to the honourable Curatorium to have the university treasury make up, for the present widows, as many of them as are then still alive, the amount by which their pensions fall short of 250 Rthl. during the period in which the above regulation causes the pension to fall below 250 Rthl.
\end{quote}
\noindent\hspace*{\parindent}Such a request could be supported by readily apparent grounds. It is also easy to see that such assistance could impose no great burden on the university treasury. It would not even be impossible for the need for assistance \textit{never} to arise. But even if it were needed soon, because a new widow joined their number before one of the present widows had left it, the supplement for each widow could still be expected to amount to only a few thalers for a considerable time. In the more distant future, for example after 30 years, if we assume by way of example one of the most extreme cases, namely that no increase in the pension had by then become permissible and that the number of widows had risen to 30, the supplement for one widow could amount to 20 Rthl.; but by then probably only a few of the present widows would still be alive.

The commission believes that it must also mention here a couple of other arrangements put forward as ways of improving the balance, without, however, wishing to make a formal proposal for them.

First, the question arose whether some improvement in the balance could be achieved by giving extraordinary professors joining in future a claim only to a lower widow's pension. They would thus form a second class of participants, from which they would automatically rise to the first on promotion to an ordinary professorship; the contributions of members of the second class, however, were to remain at their existing level of 2 Louis d'or. The present extraordinary professors were to be given the choice of remaining in the first class or moving to the second, but in the latter case a return to the first class would be inadmissible as long as they remained extraordinary professors.

To assess this idea properly, it must first be borne well in mind that, in the main plan set out above on p. [160], the reduction of the deficit from 17520 Rthl. to 2963 Rthl. depends essentially on the assumption that the contributions of \textit{all} members, present and future alike, are raised from 2 to 3 Louis d'or. That main plan therefore cannot stand at all if this assumption suffers a substantial shortfall without full and secure compensation being provided elsewhere. This consideration itself rules out as inadmissible a way of setting the pension of the second class lower which would otherwise commend itself by its simplicity: namely, to allow widows of the second class only the fixed part of the pension, 200 Rthl., or to place them entirely on the same footing as Professor G.'s widow. For it is in fact easy to see that, very probably, \textit{the fund} would thereby receive no compensation at all for that loss of contributions, and only the widows of the first class would be somewhat better off, since the supplementary dividend would be shared among fewer participants.

